\subsubsection{Contraction of Discrete-Time Nonlinear Systems}

\citet{tran2019convergence} study the nonlinear discrete-time system
$\mathbf{s}_{t+1}=\mathcal F(t,\mathbf{s}_t)$ and distinguish convergence to
one bounded trajectory from incremental stability, which compares any two
trajectories. Their contraction condition follows the variational dynamics
through the Jacobian. In their time- and state-dependent metric $\Theta$, the
transformed Jacobian is

\begin{equation}
    \mathcal J_{\Theta}(t,\mathbf{s})
    =
    \Theta(t+1,\mathbf{s})
    J_{\mathcal F}(t,\mathbf{s})
    \Theta(t,\mathbf{s})^{-1},
\end{equation}

and uniform contraction requires bounded, nonsingular $\Theta$ and
$\mathcal J_{\Theta}^{\top}\mathcal J_{\Theta}\preceq(1-\mu)I$ for some
$\mu>0$. The paper proves that, for continuously differentiable dynamics,
uniform contraction is equivalent to uniform exponential incremental
stability. On a compact positively invariant domain it also implies exponential
convergence to a unique bounded trajectory. The present benchmark uses the
corresponding fixed-norm sufficient condition
$\sup_{\mathbf{s}\in D}\|J_{\mathcal F}(\mathbf{s})\|<1$: two executions
driven by the same recurrence become geometrically closer after every step.
